\chapter{Аналитическая часть}

Задача коммивояжёра заключается в поиске маршрута минимальной длины, проходящего по всем вершинам графа ровно один раз и возвращающегося в исходную вершину. Такой маршрут называется гамильтоновым циклом. В работе рассматривается неориентированный полносвязный граф, на рёбрах которого заданы затраты по времени. Дополнительно вводится ограничение на суммарную длину маршрута: допустимыми считаются только циклы с общей продолжительностью не более 80 условных дней. Общие сведения о задачах дискретной оптимизации и маршрутизации можно найти в~\cite{malyshev}.

\section{Основные понятия теории графов}

Графом называется пара \(G = (V, E)\), где \(V\) --- непустое конечное множество вершин, а \(E\) --- множество рёбер, каждое из которых соединяет две вершины из \(V\). Если каждому ребру \((i, j)\) сопоставлено числовое значение \(c_{ij}\), то говорят о взвешенном графе, а величина \(c_{ij}\) интерпретируется как стоимость или длина перехода между вершинами \(i\) и \(j\).

Граф называется неориентированным, если ребро \((i, j)\) не имеет направления, то есть переход из вершины \(i\) в вершину \(j\) эквивалентен переходу из \(j\) в \(i\), и вес ребра удовлетворяет равенству \(c_{ij} = c_{ji}\). Граф называется полносвязным (полным), если каждая пара различных вершин соединена ребром.

Маршрутом в графе называется последовательность вершин, в которой каждые две соседние вершины соединены ребром. Если маршрут начинается и заканчивается в одной и той же вершине и при этом каждая вершина графа встречается в нём ровно один раз, такой маршрут называется гамильтоновым циклом. В рассматриваемой работе используется взвешенный неориентированный полносвязный граф, на рёбрах которого заданы затраты по времени, а искомый маршрут представляет собой гамильтонов цикл с ограничением на суммарную длину.


\section{Метод полного перебора}

Метод полного перебора предполагает просмотр всех возможных перестановок вершин графа. В рассматриваемой постановке фиксируется начальная вершина (например, вершина с номером ноль), после чего перебираются все перестановки остальных вершин. Для каждой перестановки строится гамильтонов цикл с возвратом в исходную вершину и вычисляется его длина. Если длина превышает заданное ограничение, маршрут отбрасывается. Подробное изложение классических методов перебора и оценки их трудоёмкости приведено, в частности, в~\cite{malyshev}.

Преимущества метода:
\begin{itemize}
	\item гарантированное нахождение оптимального по длине маршрута среди всех допустимых циклов;
	\item простота идеи и реализации;
\end{itemize}

Недостаток метода:
\begin{itemize}
	\item невозможность применения для графов большого размера из-за экспоненциального роста времени выполнения;
\end{itemize}

Метод полного перебора используется в работе как эталон для оценки качества приближённых решений, получаемых вероятностными методами.

\section{Муравьиный алгоритм}

Муравьиный алгоритм относится к стохастическим методам оптимизации и основан на моделировании поведения муравьиной колонии при поиске кратчайших путей между гнездом и источниками пищи. Общее изложение идей муравьиных алгоритмов и их применения к задачам маршрутизации приведено, например, в~\cite{dorigo}. В контексте задачи коммивояжёра каждый муравей строит маршрут, проходящий по всем вершинам графа и возвращающийся в исходную вершину. Выбор следующей вершины при построении маршрута выполняется с учётом двух факторов:
\begin{itemize}
	\item количества феромона на ребре, характеризующего предпочтительность данного перехода;
	\item локальной эвристики, связанной с длиной ребра (чем меньше длина, тем выше привлекательность перехода);
\end{itemize}

Вероятность выбора очередной вершины определяется формулой~(\ref{eq:prob}):
\begin{equation}
	p_{ij} \propto \tau_{ij}^{\alpha} \eta_{ij}^{\beta},
	\label{eq:prob}
\end{equation}

где \(\tau_{ij}\) --- количество феромона на ребре \((i,j)\), \(\eta_{ij} = 1 / c_{ij}\) --- эвристическая оценка, \(\alpha\) и \(\beta\) --- настраиваемые параметры.

В работе используется муравьиный алгоритм без элитных муравьёв, при этом вводится нижняя граница \(\tau_{\min}\) на значение феромона, что не позволяет вероятности перехода по ребру к непосещённой вершине, которая может стать нулевой~(\ref{eq:prob}). После построения маршрутов всеми муравьями выполняется обновление феромона: часть существующего количества испаряется, а на рёбра допустимых маршрутов (с длиной не более 80) добавляется феромон, величина которого обратно пропорциональна длине найденного цикла. Маршруты, нарушающие ограничение по длине, в обновлении феромона не участвуют.

Преимущества муравьиного алгоритма:
\begin{itemize}
	\item возможность получения приближённых решений для графов большей размерности по сравнению с полным перебором;
	\item наличие механизма самообучения за счёт накопления феромона на перспективных маршрутах;
	\item естественный параллелизм при моделировании нескольких муравьёв;
\end{itemize}

Недостатки муравьиного алгоритма:
\begin{itemize}
	\item отсутствие гарантии нахождения оптимального решения;
	\item чувствительность качества и скорости сходимости к выбору параметров алгоритма (\(\alpha\), \(\beta\), \(\rho\));
	\item необходимость проведения серии запусков для получения статистически устойчивых оценок;
\end{itemize}

В данной работе предусмотрена параметризация муравьиного алгоритма по трём основным параметрам с целью анализа влияния их значений на качество итоговых маршрутов и выбора рекомендованных настроек для рассматриваемого варианта задачи коммивояжёра.

\section{Вывод}

В аналитической части сформулирован рассматриваемый вариант задачи коммивояжёра, описаны два подхода к её решению --- метод полного перебора и метод на основе муравьиного алгоритма с учётом ограничений варианта. Определены их основные преимущества и недостатки, а также обозначена необходимость параметризации муравьиного алгоритма для последующего исследования качества получаемых решений.
